% =========================================================
% DRAFT - Frequency-Constrained Unit Commitment with
% Multi-Source Virtual Inertia from Inverter-Based Resources
% Author  : Nathaniel Benelisha (UGM)
% Version : v0.1 - Draft (30 August 2026)
% Target  : IEEE Access / Energies (MDPI)
% STATUS  : \RESULT{X} = placeholder, replace after simulation
% =========================================================

\documentclass[journal]{IEEEtran}
\usepackage{amsmath,amssymb}
\usepackage{graphicx}
\usepackage{booktabs}
\usepackage{cite}
\usepackage{xcolor}
\newcommand{\RESULT}[1]{\textcolor{red}{\textbf{[#1]}}}

\begin{document}

\title{Frequency-Constrained Unit Commitment with\\
Multi-Source Virtual Inertia Provision from\\
Inverter-Based Resources}

\author{Nathaniel~Benelisha~and~Lesnanto~Multa~Putranto,~\IEEEmembership{Ph.D., IPM., SMIEEE}
\thanks{Department of Electrical Engineering and Information Technology,
Universitas Gadjah Mada, Yogyakarta 55281, Indonesia.
E-mail: nathanielbenelisha@mail.ugm.ac.id}}
\maketitle

% =========================================================
\begin{abstract}
The rapid proliferation of Inverter-Based Resources
(IBR)---including solar photovoltaic (PV), wind turbines,
and Battery Energy Storage Systems (BESS)---reduces natural
rotational inertia, causing higher Rate of Change of
Frequency (RoCoF) and deeper frequency nadirs after N-1
contingencies. This paper proposes an enhanced
Frequency-Constrained Unit Commitment (FCUC) that
co-optimizes conventional generator scheduling and virtual
inertia from BESS, solar PV, and wind turbines simultaneously
within a Mixed-Integer Linear Programming (MILP) framework.
A two-stage linearization converts the nonlinear frequency
nadir constraint into a tractable linear form. RoCoF, nadir,
and quasi-steady-state (QSS) constraints are enforced under
worst-case N-1 contingency. Case studies on a modified IEEE
10-unit system (20\%--80\% IBR penetration) show the
proposed model achieves \RESULT{COST\_SAVING}\% lower cost
versus synchronous-generator-only inertia, with implications
for Indonesia's RUPTL 2025--2034 renewable targets.
\end{abstract}

\begin{IEEEkeywords}
FCUC, Virtual Inertia, IBR, BESS, RoCoF, Frequency Nadir, MILP.
\end{IEEEkeywords}

% =========================================================
\section{Introduction}

\subsection{Motivation}
Indonesia's RUPTL 2025--2034 targets 69.5~GW of new
generation with 61\% from renewables~\cite{PLN2025}.
Increasing IBR penetration reduces system inertia, causing
higher RoCoF~\cite{Tielens2016} and deeper nadirs below the
Indonesian UFLS threshold of 49.0~Hz~\cite{ESDM2007}.
High-profile incidents (South Australia 2016, Great Britain
2019) demonstrate these risks~\cite{AEMO2017}.

\subsection{Literature Review}
Classical MILP UC minimizes cost but ignores frequency
dynamics~\cite{Padhy2004}. Wen \textit{et al.}~\cite{Wen2016}
first embedded analytical RoCoF and nadir constraints with
BESS (foundational FCUC paper). Chu \textit{et al.}~\cite{Chu2020}
co-optimized synthetic inertia from wind in a stochastic UC.
Shi \textit{et al.}~\cite{Shi2025} linearized the nadir
constraint via two-stage regression. Indonesian studies on
Sulawesi~\cite{Sulawesi2024} and JAMALI~\cite{JAMALI2025}
provide local benchmarks.

\subsection{Contributions}
\begin{enumerate}
  \item Multi-source VI FCUC: unified MILP with BESS,
    solar PV, and wind VI simultaneously.
  \item Extended nadir linearization incorporating IBR VI.
  \item Cost-security trade-off across 20\%--80\% IBR.
  \item Indonesia RUPTL 2025--2034 alignment.
\end{enumerate}

% =========================================================
\section{System Modeling}
\label{sec:model}

\subsection{Conventional Generator}
Quadratic fuel cost:
\begin{equation}
  C_i^{gen}(P_{i,t}) = a_i P_{i,t}^2 + b_i P_{i,t} + c_i
  \label{eq:cost}
\end{equation}
Inertia constant {SG,i}$ [MWs/MVA] (typical: 2--9~s).

\subsection{BESS Model}
SoC dynamics ($\eta_c$, $\eta_d$ = efficiencies):
\begin{equation}
  E_t = E_{t-1} + \eta_c P_{c,t}\Delta t
          - \frac{1}{\eta_d}P_{d,t}\Delta t
  \label{eq:soc}
\end{equation}
Virtual inertia via grid-forming control:
\begin{equation}
  H_{VI,BESS,t} = K_{VI}^{BESS}\frac{P_{BESS,t}^{avail}}{S_{sys}}
  \label{eq:hbess}
\end{equation}

\subsection{Solar PV (PLTS)}
Synthetic inertia loop:
\begin{equation}
  H_{VI,PV,t} = K_{VI}^{PV}\frac{P_{PV,t}}{S_{sys}}
  \label{eq:hpv}
\end{equation}
Deloading: {PV,t} \leq (1-\delta_{PV})P_{PV,t}^{avail}$.

\subsection{Wind Turbine (PLTB)}
Rotor kinetic energy extraction~\cite{Chu2020}:
\begin{equation}
  H_{VI,WT,t} = K_{VI}^{WT}\frac{P_{WT,t}}{S_{sys}}
  \label{eq:hwt}
\end{equation}
Deloading: {WT,t} \leq (1-\delta_{WT})P_{WT,t}^{avail}$.

\subsection{Total System Inertia}
\begin{equation}
  H_{sys,t} = \!\sum_{i\in\mathcal{G}}\!
    \frac{H_{SG,i}S_i}{S_{sys}}u_{i,t}
    + H_{VI,BESS,t} + H_{VI,PV,t} + H_{VI,WT,t}
  \label{eq:hsys}
\end{equation}

\subsection{System Frequency Response}
Post-contingency swing equation:
\begin{equation}
  2H_{sys,t}\frac{d(\Delta f)}{dt}
    = -\Delta P_{loss} + \Delta P_{gov}
    + \Delta P_{VI} - D\Delta f
  \label{eq:swing}
\end{equation}
Key metrics:
\begin{align}
  \text{RoCoF} &=
    \frac{f_0\,\Delta P_{loss}}{2H_{sys,t}S_{sys}}
  \label{eq:rocof}\\
  \Delta f_{QSS} &=
    \frac{f_0\,\Delta P_{loss}}{DS_{sys}+R_{sys,t}}
  \label{eq:qss}
\end{align}
where {sys,t}=\sum_i R_i u_{i,t}+R_{VI,t}$ is droop.

% =========================================================
\section{FCUC Formulation}
\label{sec:formulation}

\subsection{Objective Function}
\begin{equation}
  \min\!\sum_{t=1}^T\!\Bigg[\sum_{i\in\mathcal{G}}
    \!\Big(C_i^{gen}+C_i^{SU}v_{i,t}+C_i^{SD}w_{i,t}\Big)
    + c_{deg}(P_{d,t}+P_{c,t})\Delta t\Bigg]
  \label{eq:obj}
\end{equation}

\subsection{Standard UC Constraints}
Power balance, generation limits, ramp rates, min up/down
times, logical consistency, and spinning reserve follow
standard MILP UC formulations~\cite{Padhy2004}.

\subsection{BESS Constraints}
SoC~\eqref{eq:soc}, bounds ^{min}\!\leq E_t\!\leq E^{max}$,
VI headroom \!\geq E^{min}+E^{VI}$, power limits,
charge/discharge exclusivity, daily neutrality =E_0$.

\subsection{IBR Constraints}
Output within available capacity; VI limited by
\eqref{eq:hpv}--\eqref{eq:hwt} and available generation.

\subsection{Frequency Security Constraints}
N-1 loss of largest committed unit (linearized via Big-M):
\begin{equation}
  \Delta P_{loss,t}\!\geq P_{i,t}\!-M(1-\phi_{i,t}),\;
  \sum_i\phi_{i,t}=1
  \label{eq:bigm}
\end{equation}

\textbf{RoCoF} (from \eqref{eq:rocof}):
\begin{equation}
  H_{sys,t} \geq
    \frac{f_0\,\Delta P_{loss,t}}
         {2\,\text{RoCoF}^{max}\,S_{sys}}
  \label{eq:rocof_con}
\end{equation}

\textbf{Nadir (two-stage linearization,}
extended from~\cite{Shi2025}\textbf{):}
\begin{equation}
  \alpha_0+\alpha_1 H_{sys,t}+\alpha_2 R_{sys,t}
  +\alpha_3\Delta P_{loss,t}
  \leq \Delta f^{max}
  \label{eq:nadir_con}
\end{equation}
Coefficients from 5000-sample SFR regression.

\textbf{QSS} (from \eqref{eq:qss}):
\begin{equation}
  R_{sys,t} \geq
    \frac{f_0\,\Delta P_{loss,t}}
         {\Delta f_{QSS}^{max}S_{sys}} - DS_{sys}
\end{equation}
\textbf{Min inertia:} {sys,t}\geq H_{sys}^{min}$.

% =========================================================
\section{Case Study}
\label{sec:casestudy}

\subsection{Test System}
Modified IEEE 10-unit system~\cite{Kazarlis1996};
=24$~h, $\Delta t=1$~h. Generator parameters in
Table~\ref{tab:gen}; IBR/BESS in Table~\ref{tab:ibr};
frequency limits in Table~\ref{tab:freq}.

\begin{table}[!t]\renewcommand{\arraystretch}{1.2}
\caption{Generator Parameters~\cite{Kazarlis1996}}
\label{tab:gen}\centering\footnotesize
\begin{tabular}{ccccccc}
\toprule
Unit&^{min}$&^{max}$&$&$&$&{SG}$\\
&[MW]&[MW]&[\$/MW2]&[\$/MWh]&[\$/h]&[s]\\
\midrule
G1 &150&455&0.00048&16.19&1000&5.0\\
G2 &150&455&0.00031&17.26& 970&5.0\\
G3 & 20&130&0.00200&16.60& 700&4.0\\
G4 & 20&130&0.00211&16.50& 680&4.0\\
G5 & 25&162&0.00398&19.70& 450&4.0\\
G6 & 20& 80&0.00712&22.26& 370&3.5\\
G7 & 25& 85&0.00898&27.74& 480&3.5\\
G8 & 10& 55&0.00413&25.92& 660&3.0\\
G9 & 10& 55&0.00222&27.27& 665&3.0\\
G10& 10& 55&0.00173&27.79& 670&3.0\\
\bottomrule
\end{tabular}
\end{table}

\begin{table}[!t]\renewcommand{\arraystretch}{1.2}
\caption{IBR/BESS Parameters}
\label{tab:ibr}\centering\footnotesize
\begin{tabular}{lrc}
\toprule
Parameter&Value&Unit\\\midrule
BESS power/energy&200/400&MW/MWh\\
SoC limits&10\%--90\%&---\\
$\eta_c$, $\eta_d$&0.95, 0.95&---\\
{deg}$&5.0&\$/MWh\\
{VI}^{BESS}$&10&---\\
Wind capacity&300&MW\\
Solar PV capacity&200&MW\\
{VI}^{WT}$/{VI}^{PV}$&3.5/2.5&---\\
{sys}$, $&3000, 1.0&MVA, p.u.\\
\bottomrule
\end{tabular}
\end{table}

\begin{table}[!t]\renewcommand{\arraystretch}{1.2}
\caption{Frequency Security Limits}
\label{tab:freq}\centering\footnotesize
\begin{tabular}{lrc}
\toprule
Limit&Value&Unit\\\midrule
$&50&Hz\\
RoCoF{max}$&0.5&Hz/s\\
{nadir}^{min}$&49.0&Hz\\
{QSS}^{min}$&49.5&Hz\\
{sys}^{min}$&3.0&MWs/MVA\\
\bottomrule
\end{tabular}
\end{table}

Demand profile: Kazarlis benchmark (peak 1500~MW),
spinning reserve = 10\% of hourly load.

\subsection{Scenarios}
S0: Classical UC (no freq.\ constraints, 0\% IBR).
S1: FCUC, SG inertia only (0\% IBR).
S2: FCUC + BESS VI (40\% IBR).
S3: FCUC + BESS + Wind VI (40\% IBR).
S4: Proposed---BESS+Wind+Solar VI (40\% IBR).
S5: Proposed at 60\% IBR.
S6: Proposed at 80\% IBR.

% =========================================================
\section{Results and Discussion}
\label{sec:results}

\subsection{Nadir Linearization Validation}
The linearized nadir~\eqref{eq:nadir_con} achieves
MAE~=~\RESULT{MAE}~Hz, max error~=~\RESULT{MAXERR}~Hz,
^2$~=~\RESULT{R2} over \RESULT{NPTS} test points.

\subsection{S0 vs.\ S1: Classical UC vs.\ FCUC}
S0 costs \RESULT{S0COST}~\$/day but produces
RoCoF~=~\RESULT{S0ROCOF}~Hz/s (violates in
\RESULT{S0VIO}~periods) and nadir~=~\RESULT{S0NAD}~Hz.
S1 (FCUC): cost \RESULT{S1COST}~\$/day
(+\RESULT{S1PREM}\%), all constraints satisfied.

\subsection{S2--S4: Multi-Source VI at 40\% IBR}
\begin{table}[!t]\renewcommand{\arraystretch}{1.2}
\caption{Performance at 40\% IBR Penetration}
\label{tab:s234}\centering\footnotesize
\begin{tabular}{lccc}\toprule
Metric&S2&S3&S4 (Proposed)\\\midrule
Cost [\$/d]&\RESULT{S2C}&\RESULT{S3C}&\RESULT{S4C}\\
Saving vs S1 [\%]&\RESULT{S2SV}&\RESULT{S3SV}&\RESULT{S4SV}\\
Min {sys}$&\RESULT{S2H}&\RESULT{S3H}&\RESULT{S4H}\\
Max RoCoF [Hz/s]&\RESULT{S2R}&\RESULT{S3R}&\RESULT{S4R}\\
Min nadir [Hz]&\RESULT{S2N}&\RESULT{S3N}&\RESULT{S4N}\\
\bottomrule\end{tabular}\end{table}

\textit{Finding 1:} S4 reduces cost by \RESULT{S4SV}\%
vs.\ BESS-only S2 (fewer must-run thermal units needed).

\subsection{S4--S6: IBR Penetration Sweep}
\begin{table}[!t]\renewcommand{\arraystretch}{1.2}
\caption{Results vs.\ IBR Penetration (Proposed Model)}
\label{tab:sweep}\centering\footnotesize
\begin{tabular}{lccc}\toprule
Metric&S4 (40\%)&S5 (60\%)&S6 (80\%)\\\midrule
Cost [\$/d]&\RESULT{S4C2}&\RESULT{S5C}&\RESULT{S6C}\\
SG inertia [\%]&\RESULT{S4SGI}&\RESULT{S5SGI}&\RESULT{S6SGI}\\
IBR VI total [\%]&\RESULT{S4VII}&\RESULT{S5VII}&\RESULT{S6VII}\\
-- BESS [\%]&\RESULT{S4VB}&\RESULT{S5VB}&\RESULT{S6VB}\\
-- Wind [\%]&\RESULT{S4VW}&\RESULT{S5VW}&\RESULT{S6VW}\\
-- Solar [\%]&\RESULT{S4VP}&\RESULT{S5VP}&\RESULT{S6VP}\\
\bottomrule\end{tabular}\end{table}

\textit{Finding 2:} IBR VI share grows from
\RESULT{S4VII}\% to \RESULT{S6VII}\% as penetration
rises from 40\% to 80\%.
\textit{Finding 3:} At 80\% IBR, only the proposed
multi-source approach finds a feasible frequency-secure
schedule.

\subsection{Economic Summary}
\begin{table}[!t]\renewcommand{\arraystretch}{1.2}
\caption{Economic and Security Summary}
\label{tab:economic}\centering\footnotesize
\begin{tabular}{lccl}\toprule
Scenario&Cost [\$/d]&vs.\ S0 [\%]&Secure\\\midrule
S0 (Classic)&\RESULT{EC0}&---&$\times$\\
S1 (SG only)&\RESULT{EC1}&+\RESULT{EP1}&$\surd$\\
S2 (BESS)   &\RESULT{EC2}&+\RESULT{EP2}&$\surd$\\
S4 (40\%)   &\RESULT{EC4}&+\RESULT{EP4}&$\surd$\\
S5 (60\%)   &\RESULT{EC5}&+\RESULT{EP5}&$\surd$\\
S6 (80\%)   &\RESULT{EC6}&+\RESULT{EP6}&$\surd$\\
\bottomrule\end{tabular}\end{table}

\textit{Finding 4:} Security premium decreases with IBR
penetration; at $\approx$\RESULT{BEVEN}\% IBR, cost of
multi-source FCUC approaches classical UC.

% =========================================================
\section{Conclusion}
\label{sec:conclusion}

Multi-source virtual inertia FCUC is essential for
high-IBR power systems. Key conclusions:
(1)~Proposed model achieves \RESULT{FSAV}\% lower cost
than SG-only FCUC while satisfying all frequency
constraints at all tested IBR levels.
(2)~Extended two-stage nadir linearization achieves
\RESULT{FMAE}~Hz MAE and solves in \RESULT{FTIME}~s.
(3)~At 80\% IBR, multi-source VI enables feasibility
where single-source fails.
(4)~For Indonesia's RUPTL 2025--2034 (61\% renewables
by 2034), planned BESS and IBR assets can satisfy
inertia requirements without disproportionate cost.
Future work: stochastic UC; DC network constraints;
JAMALI/Sulawesi validation; joint angle+frequency stability.

\bibliographystyle{IEEEtran}
\begin{thebibliography}{20}
\bibitem{PLN2025}PLN, `RUPTL PLN 2025--2034,'' 2025.
\bibitem{Tielens2016}P.~Tielens and D.~Van~Hertem, \textit{Renew. Sustain. Energy Rev.}, 55, 999--1009, 2016.
\bibitem{ESDM2007}Permen ESDM No. 3/2007: Aturan Jaringan Jawa-Madura-Bali.
\bibitem{AEMO2017}AEMO, `Black System SA 2016 -- Final Report,'' 2017.
\bibitem{Padhy2004}N.~Padhy, \textit{IEEE Trans. Power Syst.}, 19(2), 1196--1205, 2004.
\bibitem{Wen2016}Y.~Wen \textit{et al.}, \textit{IEEE Trans. Power Syst.}, 31(6), 5115--5125, 2016.
\bibitem{Chu2020}Z.~Chu \textit{et al.}, \textit{IEEE Trans. Power Syst.}, 35(5), 4056--4066, 2020.
\bibitem{Shi2025}Q.~Shi \textit{et al.}, \textit{IEEE Trans. Power Syst.}, 40(4), 3584--3587, 2025.
\bibitem{Sulawesi2024}[Authors], Unit Commitment with PFR, Southern Sulawesi, 2024.
\bibitem{JAMALI2025}[Authors], Unit Commitment with FFR, JAMALI, ICT-PEP, 2025.
\bibitem{Kazarlis1996}S.~A.~Kazarlis \textit{et al.}, \textit{IEEE Trans. Power Syst.}, 11(1), 83--92, 1996.
\bibitem{Kundur1994}P.~Kundur, \textit{Power System Stability and Control}, McGraw-Hill, 1994.
\bibitem{Anderson1990}P.~M.~Anderson and M.~Mirheydar, \textit{IEEE Trans. Power Syst.}, 5(3), 720--729, 1990.
\bibitem{Asiedu2025}S.~T.~Asiedu \textit{et al.}, \textit{IEEE Access}, 13, 2025.
\end{thebibliography}

\end{document}